\section[经典 Langlands]{经典 Langlands（6 月 3 日）}

\srcpage{85}

说明：我（Arun）对数论了解得不够好，无法忠实记下今天的课，而今天讲得很快。我已尽力，但会格外留意笔误、数学错误与失误。谢谢理解。

回忆几何 Satake 定理讲的是自守侧 TFT 中的线算子：$(\cA_G(S^2),*)\cong\Rep_{G^\vee}$。这些给出 't Hooft 线算子，这种算子的荷（charge）是一个 $G^\vee$-表示（或它在表示环中的类）。还有 't Hooft 环算子：四维流形 $M$ 中环的链环（link）是 $S^2\times S^1$，所以环算子是向量空间 $\cA_G(S^2\times S^1)$。

\srcpage{86}

正如上次讨论的，与 $S^1$ 作叉积是某种维数或迹的推广；这里我们取 $K_0$，于是 $\cA_G(S^2\times S^1)\cong K_0(\Rep_{G^\vee})$。这就是表示环。

给定带纽结 $K\subset Y$ 的三维流形 $Y$，我们可以用 $K_0(\Rep_{G^\vee})$ 的元素 $V$ 给 $K$ 贴标签。这定义了一个算子：把纽结放进 $Y\times[0,1]$ 的 $1/2$ 处。$K_0(\Rep_{G^\vee})$ 也标注 $B_{G^\vee}$ 中的 Wilson 环；它们是更简单的算子，通过研究平坦 $G^\vee$-联络得到：给定纽结 $K\subset Y$，我们取表示 $V$ 中沿 $K$ 的和乐（holonomy）的迹。这定义了 $\LocSys_{G^\vee}(Y)$ 上的函数。电磁对偶应当把 A 侧与 B 侧的向量空间等同起来；一侧我们有相对复杂的算子，而这一侧我们有特征簇（联络空间）上最简单的一类函数。这两类相当不同的算子由同一个集合 $K_0(\Rep_{G^\vee})$ 标注，而电磁对偶应当把它们互换。一般地，这个故事很神秘；不借助猜想，我们对较难的 A 侧那些向量空间理解得并不好，而那是作弊。

另一方面，在经典 Langlands 纲领中，事情稍微模糊一些。三维流形 $Y$ 的类比物是有限域上的曲线，或数域。最重要的特征是起作用的算子代数：$Y$ 中的纽结被换成 $C$ 上的点或 $\cO_F$ 中的素元。在几何设定中，算子向量空间来自几何 Satake 范畴的去范畴化；而在算术设定中，$K_0(\Rep_{G^\vee})$ 以更微妙的方式去范畴化地出现：取 Frobenius 在球面 Hecke 范畴上的迹，而不是恒等的迹。具体地，$K_0(\Rep_{G^\vee})$ 出现为 Frobenius 在 $\mathrm{Shv}(G(\cO)\backslash G(K)/G(\cO))$ 上的迹。这里 $K=\F_q((t))$，$\cO=\F_q[[t]]$。朴素地去范畴化会说：考虑这个双陪集空间上的函数；经典 Satake 定理把函数空间与 $K_0(\Rep_{G^\vee})$ 等同起来。

（TODO：漏了一堆东西，抱歉。）

经典 Satake 定理由 Satake 与 Langlands 在几何 Satake 定理被表述之前几十年就证明了；我们用这种时代错置的方式呈现它，是为了说明它如何嵌入我们一直在讲的故事。

同样重要的问题是：为什么我们经典地研究这个 Hecke 代数？设 $C$ 是有限域上的曲线；则我们把 $\BunG(C)$ 的点与
\[
G(F)\backslash G(\A_F)/G(\cO_\A)
\]
等同起来，其中 $\A$ 是 adele 的（TODO：$F$ 是全局域（TODO：？？），而 $\cO_\A$ 是 adele 的整数环（？？）构成的环。TODO：抱歉各位！）\yizhu{原文此处作 acdeles，后文又作 Laffourge，均系笔误，已按 adeles、Lafforgue 译出。}{B（原文笔误）}我们还把它与限制乘积的双商
\[
G(F)\Big\backslash\prod_x{}' G(K_x)/G(\cO_x)
\]
等同起来。

于是，若取 $\BunG(C)$ 上的函数，则给定 $C$ 中一点 $x$ 的选择，球面 Hecke 代数就作用；一般地，所有 $x\in C$ 上 $K_0(\Rep_{G^\vee})$ 的限制张量积作用在 $\BunG(C)$ 上的函数空间。TODO：漏了点什么。

局部系统的空间是 $\{\Gal(F)\to G^\vee\}/G^\vee$，其中 $G^\vee$ 通过共轭作用。这是 Galois 理论与 étale 覆盖理论的等同。取 $x\in C$，我们可以问 $x$ 周围小圆周上的单值（monodromy）是什么；这定义了映射 $\LocSys_{G^\vee}(C)\to G^\vee/G^\vee$；它对应于取一个 $\Gal(F)$-表示并在 Frobenius 上求值。$G^\vee/G^\vee$ 上的函数（我们把它与 $K_0(\Rep_{G^\vee})$ 等同）映到 $\cO(\LocSys_{G^\vee}(C))$，而这就是 Wilson 线算子代数。

现在，TODO 引出一个问题：$\BunG(C)$ 上的函数空间是否在 $\LocSys_{G^\vee}(C)$ 上层化？就像我们在几何设定中看到的，A 侧是否在 B 侧上层化？这是 Lafforgue 的一个定理，叫做「自守到谱方向」。一旦你知道它是模，就可以问它是哪个模，而这是 [AGK+20b, AGK+20a, AGK+21] 的一个精确猜想：我们得到 $\omega(\LocSys_{G^\vee}(C))$ 作为 $\LocSys_{G^\vee}(C)$ 上的层。Lafforgue 的结果是我们在函数域 Langlands 纲领中所知最强的东西之一——他甚至证明了分歧情形的一个版本。

说到分歧设定，我们把它推广一下。谈到分歧，有两种层次／世代／做法：一个预备步骤，然后一个更精细的步骤。当我们说自守侧的「分歧」时，我们把 $\BunG(C)$ 换成带额外数据的丛：某个有限子集 $S\subset C$ 上的级结构（level structure）。我们对这一点保持不可知，但实践中这涉及在这些点选择标签：某种额外结构。从 adele 的观点看这看起来很自然：不是在右边对整个 $\prod G(\cO)$（好吧，其实是限制乘积）

\srcpage{87}

取商，我们只包含不在 $S$ 中的点。而在数域设定中，不是对所有素元取限制乘积的商，我们只取不在 $S$ 中的素元，并对 $p\in S$ 使用 $K_p$。

第一种朴素的分歧做法是忽略 $S$。远离它会发生什么？TODO：漏了一些东西。

无论如何，球面 Hecke 代数仍在不在 $S$ 的每一点作用：我们仍有 $G(K)/G(\cO)$，所以仍保留它。于是带此级结构的自守形式空间仍有那个巨大限制张量积的作用，但这次只在 $C\setminus S$ 的点上。这比之前小一点，那又怎样；它仍然巨大。这是一种有点傻的代数：它对 $C\setminus S$ 的拓扑一无所知。而它映到 $\cO(\LocSys_{G^\vee}(C\setminus S))$；这是一个商，而且是很好的商。它可以被解释为强迫邻近点的独立因子彼此对话。于是对谱分解，我们可以问：这个作用是否通过这个商分解？再次，在函数域设定中，这是 Lafforgue 的定理。也就是说，对任意分歧（任意级结构），自守形式在 $\LocSys_{G^\vee}(C\setminus S)$ 上层化。在无分歧点，$\Gal(F)$ 通过商 $\Z F$ 作用。

这是分歧 Langlands 对应的粗糙版本——「粗糙」也许并不恰当，因为它仍是一个非凡的定理。但我们可以进一步精化，说的是层，而不只是说它层化。这需要更好地理解分歧点处发生的事，而不只是忽略它。

刚才那是函数域设定；现在是数域。设 $F$ 是 $\Q$ 的有限扩域，$\mathrm{Aut}_{F,S,K_S}$ 是 $F$ 上自守形式的空间，依赖于级结构 $K_S$（$S$ 是 $\cO_F$ 的有限素元集）。这就是此情形下的分歧设定。你可能会问：所有 $p\notin S$ 上 $K_0(\Rep_{G^\vee})$ 的限制张量积是否作用——但这是错的！

Langlands 猜想存在一个 Langlands 群 $\mathrm{Lang}_F$，使得函数域设定的图景仍然成立，只是用 $\mathrm{Lang}_F$ 代替 Galois 群。应当有映射 $\mathrm{Lang}_F\to\Gal(F)$。某种意义上，必须有一个更大的参数空间。这还没有被定义，这也是此处难以写下 Langlands 猜想精确版本的部分原因。希望在于以 Tannakian 方式从自守形式空间提取出 Langlands 群。

在自守形式空间内部，有更小的空间：代数自守形式（algebraic automorphic form）与上同调自守形式（cohomological automorphic form），它们比通常的自守形式空间更易把握，并且捕捉到你因 Langlands 理由而关心的大部分东西。它们去掉了我们一般对自守形式所需的 $L^2$ 分析，代之以代数：例如，上同调自守形式取算术局部对称空间的上同调，这更容易。你可能想要一个联系这些不同定义的 Hodge 定理。

代数自守形式与 Galois 表示之间应当有 Langlands 式的关系；这就是 Clozel 猜想。这些 Galois 表示的系数在这里是重要的细节；你应当真的对所有 $\ell\neq p$ 的素元 $\ell$ 都有相容的 $\ell$-进 Galois 表示族。这是人们称为动机（motive）的东西的一个版本——事实上，Clozel 猜想正是用动机来陈述的：某类动机式 Galois 表示对应于自守形式。

TODO：等等，椭圆曲线是怎么出现的？

无论如何，Langlands 纲领更传统地是用模形式的 L-函数引入的，这为进入分析侧提供了方便的门。TODO：有些东西我没懂，就是「Artin 猜想」及相关的东西，还有更多。抱歉各位。TODO。

Galois 群到 $\mathrm{GL}_n(k)$ 的 Artin 表示来自零维簇。这已经很深刻，但与我们试图思考的东西并没有真正互动。无论如何，你需要（TODO：发生了什么？）考虑 $\ell$-进局部系统，所以你看的不是带 $\C$ 点的群，而是带 $\overline{\Q}_\ell$ 点的群。作为域它们当然同构，但拓扑不同。这里有些有趣的复杂性：你可以取 $\ell=p$（？？？），而那种情形必须分开处理。

若你做整体 Langlands，试图一次处理所有素元，而不是局部情形中某个特定的 $p$，就不太清楚如何选取这样的 $\ell$。这引向相容系统族（compatible families of systems）的想法：表述它意味着什么。TODO：没懂。这来自代数几何：设 $X$ 是域上的概形，你想考虑 $X$ 的上同调。一般不能用有理系数做这件事，但可以用 $\Q_\ell$ 系数（TODO：我们选了一个素元 $p$，使得 $X$ 的特征是 $p$？$\ell\neq p$）。于是 Galois 群作用在这个上同调上。我们希望「相容的系统族」有类似的感觉。事实上，我们希望这类对象是动机：某个假想的 $\Q$-线性范畴中的对象，而取 étale 上同调应当允许从一个动机中提取出 $\Q_\ell$-向量空间。动机是

\srcpage{88}

一个用来指一组 $\Q_\ell$-向量空间的词，这些向量空间全都是把不同的上同调理论作用于 $X$ 得到的，以及诸如此类。回到 Langlands 纲领，我们希望／期待代数自守形式被等同于动机式 Galois 表示。

还有一条更经典、更具体的路来做同样的事，从相同的数据出发。最终会很相似，但风味不同。我们有整体 Hecke 代数作用在自守形式空间上，我们做朴素的事，只看无分歧点上的限制乘积。谱分解是对角化的花哨版本，所以我们可能首先要找的是 Hecke 特征函数，即同时是所有 Hecke 算子的特征函数。还有增长条件，分析味道的，我们将忽略。

无论如何，如果你有这样的特征函数 $f$，如何计算其特征值？回忆 Hecke 代数是 $\cO(G^\vee/G^\vee)$ 的许多副本的巨大张量积。有一个比 TODO 更优雅的办法：取一个表示 $V\colon G^\vee\to\mathrm{GL}_n$，于是得到从 $\mathrm{GL}_n/\mathrm{GL}_n$ 上函数到 $G^\vee/G^\vee$ 上函数的映射。这就是我们谈特征标时所做的。每个人都有自己的办法来测量矩阵，那就是特征多项式；对 $G^\vee$，一般不选 $V$ 就没有好办法做到这一点。给定矩阵 $F$，特征多项式是 $\det(1-tF)$。一种花哨的说法是：考虑 $\Lambda^\bullet\C^n$；取 $F$ 在这个表示上的分次迹，这就是特征多项式。现在把 $\C^n$ 换成 $V$，我们就可以对 $G^\vee$ 的元素做这件事。数论学家有时习惯用 $1/\det(1-tF)$ 而不是特征多项式，尽管两者含有相同的信息。这个变体是 $F$ 作用在 $\mathrm{Sym}^\bullet V$ 上的分次迹。这是几何级数公式 $1/(1-t)=1+t+t^2+\cdots$ 的一个版本，也是 Koszul 对偶的一个版本；Koszul 对偶互换外代数与对称代数。

于是自守形式的 L-函数是 Hecke 特征值的特征多项式的倒数。\footnote{我很抱歉用了这么糟糕的记号，但我是在现场做这件事，一时想不出更好的！}
也就是说，它是
\[
\prod_{p\notin S}\frac{1}{\det(1-p^{-s}F_p)},
\]
其中 $F_p\in G^\vee/G^\vee$ 是 TODO。

对，所以 L-函数多多少少是花哨的特征多项式：疯狂的生成函数。TODO。你可以对 Galois 表示的 L-函数做类似的事：
\[
L(\rho,V,s)=\prod_{p\notin S}\frac{1}{\det(1-F_\rho p^{-s})}.
\]
（TODO：$p$ 与 $\rho$ 没弄清楚，当心／抱歉！）

联系自守形式与 Galois 表示的 Langlands 对应，包含它们的 L-函数也应当匹配的希望。L-函数的某些值可以在场论上解释为某个四维流形上某个 TFT 的配分函数。多数时候这样做会得到无穷大，但有办法处理。有趣的是，L-函数设定中的分析困难与物理设定中的困难看起来如此相似。
